% generado por el cuaderno homónimo; no editar
AA & 3 & \texttt{gc} (gulp call), \texttt{hm} (hoot call), \texttt{sc} (squeak call) \\
AC & 6 & \texttt{chc} (chitter call), \texttt{gc} (growl call), \texttt{cc} (contact call), \texttt{whc} (whinnie call), \texttt{sc} (squeak call), \texttt{bc} (bark call) \\
AS & 2 & \texttt{hc} (howl call), \texttt{bc} (bark call) \\
CC & 1 & \texttt{cc} (contact call) \\
LW & 10 & \texttt{cc} (contact call), \texttt{cs} (contact syllable), \texttt{aa} (aerial alarm call), \texttt{ta} (terrestrial alarm call), \texttt{sqc} (squeal call), \texttt{phc} (phee call), \texttt{tr} (trino call), \texttt{tf} (trino fast call), \texttt{tt} (trino transition call), \texttt{vc} (visual contact call) \\
PT & 2 & \texttt{dc} (duet call), \texttt{ac} (alarm call) \\
SB & 5 & \texttt{pcc} (peep contact call), \texttt{ppc} (play peep call), \texttt{lpc} (long peep call), \texttt{spc} (spit call), \texttt{sc} (shriek call) \\
SM & 8 & \texttt{cc} (contact call), \texttt{acc} (aggressive contact call), \texttt{sc} (squeal call), \texttt{pc} (purr call), \texttt{whc} (whistle call), \texttt{hic} (hip call), \texttt{fc} (food call), \texttt{fs} (food syllable)
